% !TEX TS-program = xelatex
% !TEX encoding = UTF-8

% ===================================================================
% SOTOPIA: Interactive Evaluation for Social Intelligence in Language Agents
% 中文翻译版 - 由 AI 学术翻译专家完成
% 编译命令: xelatex main.tex (需要两次以生成完整目录)
% ===================================================================

\documentclass[11pt,a4paper]{article}

% ===== 中文支持 =====
\usepackage[UTF8]{ctex}

% ===== 数学符号与公式 =====
\usepackage{amsmath,amssymb,amsfonts,amsthm}

% ===== 图形与表格 =====
\usepackage{graphicx}
\usepackage{booktabs}
\usepackage{tabularx}
\usepackage{multirow}
\usepackage{multicol}
\usepackage{array}
\usepackage{longtable}

% ===== 页面布局 =====
\usepackage[a4paper, margin=2.5cm]{geometry}

% ===== 颜色与超链接 =====
\usepackage{xcolor}
\usepackage[colorlinks=true, linkcolor=blue, citecolor=blue, urlcolor=blue, breaklinks=true]{hyperref}

% ===== 代码片段 =====
\usepackage{listings}

% ===== 章节标题微调 =====
\usepackage{titlesec}
\titleformat{\section}{\Large\bfseries\color{black!85}}{\thesection}{1em}{}
\titleformat{\subsection}{\large\bfseries\color{black!80}}{\thesubsection}{1em}{}
\titleformat{\subsubsection}{\normalsize\bfseries\color{black!75}}{\thesubsubsection}{1em}{}

% ===== 页眉页脚 =====
\usepackage{fancyhdr}
\pagestyle{fancy}
\fancyhf{}
\fancyhead[L]{\small SOTOPIA: 语言智能体的社交智能交互式评估}
\fancyhead[R]{\small ICLR 2024 中文版}
\fancyfoot[C]{\thepage}
\renewcommand{\headrulewidth}{0.4pt}

% ===== 标题信息 =====
\title{
    \vspace{-0.5cm}
    {\Large\bfseries SOTOPIA：面向语言智能体的社交智能交互式评估}\\[0.3cm]
    {\large （SOTOPIA: Interactive Evaluation for Social Intelligence in Language Agents）}\\[0.3cm]
    \large ICLR 2024 中文翻译版
}
\author{
    Xuhui Zhou$^{*}$, Hao Zhu$^{*}$, Leena Mathur, Ruohong Zhang, \\
    Zhengyang Qi, Haofei Yu, Louis-Philippe Morency, Yonatan Bisk, \\
    Daniel Fried, Graham Neubig, Maarten Sap \\[0.2cm]
    {\small 卡内基梅隆大学 语言技术研究所}\\
    {\small \texttt{https://sotopia.world}}
}
\date{\today}

\begin{document}

\maketitle
\thispagestyle{fancy}

\begin{center}
\small $^{*}$Equal contributors（共同一作）
\end{center}

\vspace{0.5cm}

% ===================================================================
\begin{abstract}
人类是社会性生物，在日常交往中追求社会目标是社会智能的关键组成部分。然而，人工智能系统在该领域的能力仍然难以捉摸。本文提出 SOTOPIA，一个用于模拟人工智能体之间复杂社交交互并评估其社交智能的开放式环境。在该环境中，智能体在多种情境下进行角色扮演并展开交互；它们相互协调、协作、交换与竞争，以达成复杂的社会目标。我们在基于大语言模型（Large Language Model, LLM）的智能体与人类之间模拟角色扮演交互，并使用名为 SOTOPIA-EVAL 的整体性评估框架对其表现进行评估。借助 SOTOPIA，我们发现这些模型在社交智能方面存在显著差异，并识别出 SOTOPIA 中一个对所有模型普遍具有挑战性的子集 SOTOPIA-hard。我们发现，在该子集上，GPT-4 的目标完成率显著低于人类，并且在社交常识推理和策略性沟通技能方面表现欠佳。这些发现表明，SOTOPIA 有望成为一个用于评估和提升人工智能体社交智能的通用研究平台。
\end{abstract}

\noindent\textbf{关键词：} 社交智能；大语言模型；多智能体交互；评估；角色扮演

\vspace{0.3cm}
\hrule
\vspace{0.3cm}

% ===================================================================
\section{引言}
% ===================================================================

人类在与他人交往中达成并平衡复杂、多元社会目标的能力，是我们作为物种所具备社交智能的重要组成部分（Kihlstrom \& Cantor, 2020; Tomasello, 2021）。即便是一个简单的社会目标，例如与朋友共享一条毯子，也需要协调自身保暖的需求与对方对个人空间的需求（图~\ref{fig:teaser}）。成功的交互要求理解他人的意图和信念（Premack \& Woodruff, 1978），同时兼顾不同乃至潜在冲突的社会规范与期望（Goffman, 1959）。

\begin{figure}[htbp]
\centering
\fbox{\parbox{0.9\textwidth}{\centering [图示说明]\textbf{图 1：SOTOPIA 概览}——一个开放式社交交互环境。在每个回合（episode）中，SOTOPIA 首先采样一个社交场景情境、社会目标和角色，然后为参与的每个智能体分配一个社会目标和角色。智能体（人工智能体或人类）在 SOTOPIA 中进行角色扮演，同时试图达成各自的目标。智能体的表现通过多维度框架 SOTOPIA-EVAL 进行评估。}}
\caption{SOTOPIA：一个开放式社交交互环境。}
\label{fig:teaser}
\end{figure}

尽管近期人工智能系统在某些情境下展现出了令人印象深刻的社交技能，但其社交智能尚缺乏稳健的评估（Shapira et al., 2023a; Ullman, 2023）。一方面，许多社交智能基准测试并非交互式的（Sap et al., 2019; Le et al., 2019; Zadeh et al., 2019b），这对于评估社交智能而言并非最优解（Mehri et al., 2022; Hoppler et al., 2022; Lee et al., 2023）。另一方面，现有的交互式评估未能充分研究多样的目标驱动行为（Zhang et al., 2018b; Park et al., 2023），或仅聚焦于特定任务（Wang et al., 2019; Padmakumar et al., 2022; FAIR et al., 2022）。

为研究动态且目标驱动的社交智能，本文提出 SOTOPIA（图~\ref{fig:teaser}），一个面向通用领域、开放式、将社交智能体置于多样化社交情境中的环境。SOTOPIA 是交互式的：在多轮模拟通信中，智能体可同时使用言语与非言语沟通以及物理动作。\footnote{动作均以文本形式表示。}它还具有多样化的任务空间：通过自动生成的情境、目标、角色、关系以及其他智能体的策略的组合，构成了一个庞大且多样的任务空间。除社会目标的完成情况外，SOTOPIA 还从多个维度评估智能体的表现。

在 SOTOPIA 中，我们创建了 90 个社交场景，涵盖合作、竞争和混合等多种社会目标，并创建了 40 个具有个性、职业、秘密、背景故事及彼此关系的角色（第~\ref{sec:env} 节），它们的笛卡尔积构成了一个庞大的任务空间。通过从该空间中采样任务，我们模拟了智能体角色扮演各自角色并基于其私密社会目标进行交互的"回合"。在该模拟中，我们不仅创建并使用了基于 LLM 的智能体，还邀请了人类参与者进行角色扮演，以研究模型与人类在社交智能方面的差异。

为评估多维度社交交互，我们不能仅考虑主要社会目标的完成情况，因为人类的动机常常需要平衡多重隐含目标，例如维护关系、保障财务、获取信息、保守秘密以及遵守社会规则。因此，我们提出 SOTOPIA-EVAL（第~\ref{sec:eval} 节），借鉴社会学、心理学和经济学领域的研究成果，从多维度对智能体进行评估。随后，我们通过人类与 GPT-4 作为评判者，将 SOTOPIA-EVAL 应用于上述模拟中的回合。我们发现 GPT-4 可以作为 SOTOPIA-EVAL 的人类判断代理，特别是在目标完成、维护财务和保持关系等维度上。

尽管更大的 LLM 通常比小模型具备更高的社交智能，但在更具挑战性的任务上（第~\ref{sec:llm} 节），它们仍难以与人类协作与竞争。它们还高度受对话伙伴的影响，存在泄露秘密和违反社会规则的风险。然而，我们也发现了一些模型能够创造性地解决问题的案例（第~\ref{sec:llm} 节）。

本文的主要贡献如下：
\begin{itemize}
    \item \textbf{(A)} 我们提出并将发布 SOTOPIA，一个用于模拟目标导向社交交互的通用领域交互式环境。SOTOPIA 具有可扩展性，可供未来研究人员研究和训练更复杂多样的社交智能体。
    \item \textbf{(B)} 我们创建了 SOTOPIA-EVAL，一个从多个社交维度分析智能体表现的多维度评估框架。
    \item \textbf{(C)} 我们通过利用 LLM 实现 SOTOPIA-EVAL 的自动化，并发现其可在部分社交维度上（尤其是目标完成）作为人类判断的代理。
    \item \textbf{(D)} 我们展示了借助 SOTOPIA 可评估模型之间的社交智能差异，以及模型与人类之间的差异。
\end{itemize}

综上，SOTOPIA 是一个新颖、富有挑战性且交互式的基准，有望成为语言智能体社交智能研究的理想测试平台和潜在的"孵化器"。

% ===================================================================
\section{SOTOPIA 交互环境}
\label{sec:env}
% ===================================================================

为应对以交互方式评估社交智能的挑战，我们寻求一个满足以下要求的环境：(1) \textbf{真实性}：能够在真实情境下评估和理解人工智能体的行为；(2) \textbf{混合效用}：人类动机往往由显式和隐式激励共同驱动，环境应能够评估智能体在多个维度上的表现；(3) \textbf{开放性}：为支持大规模模拟与评估，环境应能以过程化方式生成满足前两项要求的新任务，而无需大量人工介入。

本节将介绍 SOTOPIA，并解释为何 SOTOPIA 非常适合用于社交智能的交互式评估。任务空间包含通过人工检查的自动生成真实场景、角色与关系（第~\ref{sec:taskspace} 小节）。一个回合（episode）由扮演不同角色的智能体之间的交互组成，智能体执行动作（例如 speak("Hello Bob!")，微笑并点头，以及拨打 911）以实现从任务空间中采样的社会目标（第~\ref{sec:episode} 小节）。读者可参阅附录~\ref{app:formal} 查阅 SOTOPIA 环境的形式化定义。

\subsection{任务空间}
\label{sec:taskspace}

本文考虑涉及两个智能体的任务，但 SOTOPIA 更具普适性，可支持两个以上智能体之间的交互。SOTOPIA 中的任务由情境（scenario context）、角色及其社会目标组合而成，提供了交互的背景。每个回合由智能体之间的多轮交互组成。本文聚焦于相对较短时间跨度内单个回合中的局部一致社会目标，尽管在现实世界中，人们的社会目标会随时间持续变化。请注意，智能体对同一任务具有不同的观测：每个智能体可观测到情境、自身的社会目标以及自身的角色描述。其他智能体的社会目标不可见，其他智能体的角色描述仅部分可观测，具体取决于智能体之间的关系。

\subsubsection{任务空间的复杂度}

情境、社会目标、角色及其关系的组合可以塑造智能体的最优行为空间。考虑一个说服任务——"要求伴侣在 FaceTime 时停止发短信"。若伴侣重视从众性，则智能体达成目标的一个好方法是从社会规范的角度讨论该问题；然而，若伴侣特别善解人意，则表达主观情感可能更有效。交互伙伴的策略也会显著影响最优行为。再考虑图~\ref{fig:teaser} 所示的另一任务——"以不低于 \$3,400 的价格出售 BMW Z3"。若买家给出高报价，卖方可能希望利用买家强烈的购买意愿索要更高价格；而若买家给出过低报价，卖方可说明车辆为何更值钱或威胁离开。当交互前获取更多信息（例如个性、决策风格或职业）时，买卖双方可利用这些知识调整策略。情境、社会目标、角色、关系描述以及其他参与者策略的多样空间的笛卡尔积构成了一个庞大的任务空间，不仅提出了真实性的挑战，也为评估和提升人工智能体的社交智能提供了机会。在本小节其余部分，我们将介绍任务空间各轴的设计与生成。

\subsubsection{角色}

如上所述，角色描述的设计应包含若干影响决策的属性。我们考虑以下属性（受 Wang et al. (2019) 启发）：姓名、性别、年龄、职业、代词、人格特质（Goldberg, 1992）、道德价值观（Graham et al., 2011）、Schwartz 个人价值观（Cieciuch \& Davidov, 2012）以及决策风格（Hamilton et al., 2016），这些均通过 GPT-4（OpenAI, 2023）生成。为使对话更具背景，生成上述属性后，我们提示 GPT-4 生成秘密与公开信息。图~\ref{fig:teaser} 展示了两个角色示例。值得注意的是，尽管我们生成了多样化的角色集合，但这仅是可能角色空间的一小部分。本文的分析聚焦于以上述方式生成的 40 个角色，未来使用 SOTOPIA 的研究可轻松生成扩展的角色集。

\subsubsection{关系}

SOTOPIA 中的关系具有以下作用：(1) 情境通常对关系有约束；例如，家庭聚餐场景需要家庭关系，而在派对上寻找共同朋友的情境则不一定；(2) 不同关系会影响智能体在交互中对其他智能体描述的可见性；例如，陌生人可能不了解对方的职业，而伴侣可能知晓对方的人格。为便于 (1) 中的角色采样以及 (2) 中的交互背景控制，我们考虑五种关系类型：家人（family）、朋友（friend）、恋人（romantic）、熟人（acquaintance）和陌生人（stranger）。具体方法的局限性和潜在扩展参见附录~\ref{app:limitation}。

我们将在后续段落中讨论 (1) 的实现，对于 (2)，我们创建了一种基于规则的机制来决定角色描述的哪些部分对其他智能体可见。若两个智能体处于家人、朋友或恋人关系，它们可以互相看到对方除秘密外的所有描述。两个熟人可以互相看到对方的姓名、职业、性别代词和公开信息。两个陌生人则对彼此的描述一无所知。与角色类似，我们通过提示 GPT-4（OpenAI, 2023）根据角色池自动生成关系，并人工校验关系的一致性。

\subsubsection{情境}

我们考虑智能体同时拥有共享与私密信息的社会任务。共享信息是情境：社会交互的地点、时间以及其他共享信息，例如"一人在自家院子以 \$100 出售一把古董椅，另一人对该椅感兴趣"。私密信息是社会目标，仅对相应的智能体可见，例如"你的目标是以 \$80 买下这把椅子"仅对买家可见，而"你的目标是以 \$90 卖出这把椅子"则仅对卖家可见。然而，如上所述，情境与角色的组合并非任意，因为情境通常会对智能体施加约束。我们将此类约束称为情境约束（scenario constraints）。本文主要考虑关系约束，其决定了被采样角色之间的关系类型。与角色和关系类似，情境（包括上下文、目标和约束）均通过提示 GPT-4（OpenAI, 2023）生成。为生成覆盖多种社交交互类型的高质量情境（如图~\ref{fig:teaser} 所示），我们从先前数据集中随机采样数据（Forbes et al. 2020; Sap et al. 2019; Lewis et al. 2017; Ziems et al. 2023; He et al. 2018; 2017），将其用于提示以"启发" GPT-4。作者人工校验所有生成的情境并根据附录~\ref{app:annotate} 进行必要修改，约删除 10\% 的情境。

\subsection{SOTOPIA 回合（Episodes）}
\label{sec:episode}

在交互过程中，模型和人类被赋予社会情境、角色描述及相应的社会目标。我们将带有角色和目标的模型和人类统称为\textbf{智能体（agents）}，它们在回合中以轮流（轮询）方式执行动作，即智能体 1 先行动，然后智能体 2 行动，依次循环。在各自轮次，智能体可以选择说话、使用非言语沟通（例如拥抱或微笑）或采取物理动作（例如播放音乐），这些都是社交交互的重要组成部分（De Stefani \& De Marco, 2019）。一旦智能体从这三个离散动作类别中选择其一，便会以文本形式生成具体动作，即说什么、做什么手势等。除上述三种动作外，智能体也可选择"无动作"（none）以表示沉默或允许另一智能体完成发言，或选择"离开"（leave）以结束回合。我们将回合轮次上限设置为 20，因为我们发现人类通常可以在 20 轮内完成大部分任务。回合结束的原因要么是某智能体选择离开，要么是达到轮次上限。图~\ref{fig:teaser} 展示了一个回合示例。

% ===================================================================
\section{SOTOPIA-EVAL：整体性社交智能体评估框架}
\label{sec:eval}
% ===================================================================

为捕捉成功社交交互的复杂性，我们设计了一个受社会学、心理学和经济学文献启发的多维度评估框架。对于每个回合，智能体在交互结束时沿以下每个维度进行评分。以下段落将逐项列出 SOTOPIA 中全部七个维度，每个维度以分数范围（$[下限, 上限]$ 形式）、解释以及启发我们的文献呈现。\footnote{指标范围包含语义含义，例如，REL 中的负值表示关系恶化，正值表示关系改善。}

\begin{itemize}
    \item \textbf{目标完成度（GOAL）} $[0, 10]$：智能体达成其目标的程度。智能体的社会目标由环境定义，是其行为的主要驱动力（Weber, 1978）。
    \item \textbf{可信度（BEL）} $[0, 10]$：关注智能体的行为是否被感知为自然、真实且与角色描述一致，从而模拟可信的人类行为代理（Park et al., 2023）。具体而言，我们考虑以下标准：(1) 智能体是否以自然、真实的方式与他人交互（自然性）；(2) 智能体的动作是否与其角色特征相符，例如个性、价值观等（一致性）。
    \item \textbf{知识获取（KNO）} $[0, 10]$：捕捉智能体主动获取新信息的能力。该维度的动机在于好奇心——即渴望了解或学习的欲望是人类的基本特质（Reiss, 2004; Maslow, 1943）。具体而言，我们考虑以下标准：智能体通过交互获取了何种信息、获取的信息是否为其新信息、获取的信息对其是否重要。
    \item \textbf{保密（SEC）} $[-10, 0]$\footnote{对于 SEC 与 SOC，仅存在负值区间，因为保守秘密和遵守社会规则应被视为智能体的基线行为。}：衡量智能体（人类）对其私密信息或意图的保密需求（Reiss, 2004）。从博弈论视角看，泄露秘密通常会导致效用损失（Gilpin \& Sandholm, 2006）。然而，揭示秘密也可以成为建立信任的有力工具，从而改善关系（Jaffé \& Douneva, 2020）。在该维度，我们询问参与者希望保守何种秘密或私密意图，以及他们是否成功保守。
    \item \textbf{关系（REL）} $[-5, 5]$：捕捉人类对社会联系与归属感的基本需求（Maslow, 1943; Bénabou \& Tirole, 2006）。在该维度，我们首先询问参与者在交互前与其他智能体具有何种关系，然后评估智能体与他人的交互是否有助于保持或增强其个人关系。此外，我们还判断这些交互是否影响智能体的社会地位或声誉。
    \item \textbf{社会规则（SOC）} $[-10, 0]$：涉及规范、规章、制度安排和礼仪。我们区分两种类型的社会规则：社会规范和法律规则。法律规则包含禁止的行为及制度化强制力惩罚的潜在性；社会规范包含规范性社交规则（例如在图书馆大声说话被视为不礼貌）。
    \item \textbf{财务与物质利益（FIN）} $[-5, 5]$：涉及经典博弈论所探讨的传统经济效用（Gilpin \& Sandholm, 2006; Burns et al., 2017）。我们认为财务效用由短期货币收益（例如收入）和长期经济回报（例如工作保障、股票持有、融资机会）共同构成。
\end{itemize}

% ===================================================================
\section{研究问题与实验设置}
% ===================================================================

给定一组多样化的社交场景、目标和角色，我们模拟智能体的交互。这是我们首次能够以交互式和系统化的方式评估通用的、目标导向的社交智能体。在接下来三节中，我们将展示 SOTOPIA 如何用于研究以下问题：
\begin{itemize}
    \item \textbf{(A)} 在评估智能体社交交互时，GPT-4（OpenAI, 2023）能在多大程度上作为人类判断的代理（第~\ref{sec:gpt4-eval} 节）？
    \item \textbf{(B)} 模型之间（第~\ref{sec:llm} 节）以及模型与人类之间（第~\ref{sec:human} 节）在目标导向社交智能方面有何差异？
\end{itemize}

为研究这些问题，我们按照第~\ref{sec:env} 节中的生成流程创建了 40 个智能体、90 个关系和 90 个场景。对于每个场景，我们基于场景约束采样 5 对角色，从而得到一组 450 个任务。对于每个任务，我们通过枚举所有模型对来模拟模型间的交互。受资源限制，我们还在挑战性子集 SOTOPIA-hard（第~\ref{sec:human} 节）上模拟 GPT-4（OpenAI, 2023）与人类的交互。

具体而言，我们考虑以下模型进行比较：GPT-3.5（Ouyang et al., 2022）、GPT-4（OpenAI, 2023）、Llama-2-70b-chat（Touvron et al., 2023）以及 MPT-30b-chat（MosaicML NLP Team, 2023）。我们将智能体的温度参数设置为 1 以鼓励回复多样性，将评估者的温度参数设置为 0 以保证评估的稳定性。我们使用上述模型的固定版本以提升可复现性。\footnote{我们固定 GPT-4 版本为 gpt-4-0613，GPT-3.5 版本为 gpt-3.5-turbo-16k-0613。}

为将这些模型用作 SOTOPIA 中的智能体，在每个轮次，我们向语言模型输入情境、要扮演的角色以及交互历史以生成动作（动作类型见第~\ref{sec:episode} 小节）。本文聚焦于使用 SOTOPIA 理解社交交互，我们对 LLM 使用的提示方法与人类界面内容类似（图~\ref{fig:humaninterface}）。我们将利用新颖提示方法（如思维链（Chain-of-Thought, Wei et al., 2022）、ReAct（Yao et al., 2022））作为未来工作。

% ===================================================================
\section{GPT-4 能否评估社交交互？}
\label{sec:gpt4-eval}
% ===================================================================

本节研究以下研究问题：我们能否利用现有 LLM 自动化第~\ref{sec:eval} 节中介绍的 SOTOPIA-EVAL 评估框架？我们选择 GPT-4（OpenAI, 2023）作为本研究中的代表性模型，因为其性能优越。\footnote{在先导研究中，其他模型无法提供有意义的评估。详见附录~\ref{app:addresults}。}我们首先收集交互数据，\footnote{包括模型-人类、模型-模型和人类-人类交互。}然后要求人类基于 SOTOPIA-EVAL 中的维度对交互进行评估。\footnote{他们并不知道角色扮演者是模型还是人类。} GPT-4 被赋予与人类相同的问题集（见附录~\ref{app:eval-prompt} 与~\ref{app:annotate}），我们比较人类与 GPT-4 给出的评分。

\subsection{数据收集流程}

我们从第~\ref{sec:env} 节中随机采样了 200 个回合的子集，并在 Amazon Mechanical Turk 平台上与一组经过预先筛选的工人开展对照研究。他们被告知第~\ref{sec:eval} 节中每个维度的含义，并被展示了每个维度的高质量与低质量标注示例。他们不仅要在 11 点李克特量表（第~\ref{sec:eval} 节）上对每个智能体的每个 7 维度进行评分，还要为每个评分提供自由形式的理由。由于每个智能体的每个维度由多位人类标注者评分，我们通过对多位标注者的评分取平均来计算人类评分。标注者间一致性为中等，Randolph $\kappa$ 评分为 0.503（Randolph, 2005）。GPT-4 被赋予与人类标注者类似的任务。我们提示 GPT-4 为每个回合、智能体和维度生成一个结构化输出，包括一个整数 GPT-4 评分和理由，使用与人类相同的指令集。更多关于数据收集流程的详细信息，请参阅附录~\ref{app:annotate}。

\begin{figure}[htbp]
\centering
\fbox{\parbox{0.85\textwidth}{\centering [图示说明]\textbf{图 2：} GPT-4 评分与人类评分差异的分布。$>\!74\%$ 的 GPT-4 评分集中在人类评分 $\pm \sigma$ 范围内（$\sigma = 2.15$）。左侧白色区域大于右侧，表示当 GPT-4 与平均人类判断不一致时，其评分更可能高于人类评分。}}
\caption{GPT-4 与人类评分差异的分布。}
\label{fig:gpt4diff}
\end{figure}

\subsection{分析 GPT-4 评估与人类评估}

\begin{table}[htbp]
\centering
\small
\begin{tabular}{lcc}
\toprule
维度 & 模型 & 人类 \\
\midrule
SEC & $0.22^{**}$ & -- \\
KNO & $0.33^{**}$ & 0.19 \\
SOC & $0.33^{**}$ & $0.42^{**}$ \\
BEL & $0.45^{**}$ & $0.27^{*}$ \\
REL & $0.56^{**}$ & $0.49^{**}$ \\
FIN & $0.62^{**}$ & $0.34^{**}$ \\
GOAL & $0.71^{**}$ & $0.78^{**}$ \\
\bottomrule
\end{tabular}
\caption{GPT-4 评估与人类判断在不同维度上的皮尔逊相关系数及 $p$ 值。$^{**}: p \le 0.01$，$^{*}: p \le 0.05$。强且显著的相关系数以蓝色显示。在 GOAL 维度以及模型输出上，GPT-4 表现最佳。}
\label{tab:correlation}
\end{table}

图~\ref{fig:gpt4diff} 展示了同一维度、智能体和回合上 GPT-4 评分与人类评分的差异。我们发现，大多数（$> 74\%$）GPT-4 评分集中在人类评分 $\pm$ 一个标准差的范围内。还可以看到左侧白色区域大于右侧，这意味着当 GPT-4 与平均人类判断不一致时，其评分更可能高于而非低于人类评分。

表~\ref{tab:correlation} 将这一汇总分析按不同维度以及角色扮演者是人类还是模型进行细分。相关系数表明，当模型进行角色扮演时，GPT-4 评分与人类评分在 GOAL、FIN 和 REL 维度上具有显著且强的相关性。然而，当人类进行角色扮演时，除 GOAL 外，所有维度的相关性均显著下降。这表明 GPT-4 可在部分维度上评估社交交互，并且其在评估模型方面优于评估人类。在附录~\ref{app:addresults} 中，我们比较了单个智能体在单个回合中单个维度的平均 GPT-4 评分与人类评分范围。我们发现 GPT-4 评分通常落在大多数维度的人类评分范围内，但 SOC 和 SEC 除外，在这两个维度上 GPT-4 评分往往高于人类。

综合上述观察，我们得出结论：GPT-4 可以在一定谨慎前提下用作人类判断的代理，用于在部分维度上评估模型表现，以及在 GOAL 维度上评估人类表现。然而，我们提醒读者，LLM 作为评估者已知存在偏差和问题，包括位置偏差（Wang et al., 2023）、事实不一致（Luo et al., 2023）、偏向母语者（Liang et al., 2023）等。因此，在解读我们的结果时，应当意识到这些潜在偏差的影响。未来版本的 SOTOPIA-EVAL 可通过引入多 LLM 协作（Chan et al., 2023）以及训练更大的 LLM 评估者（Zhang et al., 2023）等近期方法进一步提升基于 LLM 的评估质量。

% ===================================================================
\section{SOTOPIA 中 LLM 间社交交互的评估}
\label{sec:llm}
% ===================================================================

我们分析模型在 SOTOPIA 中的交互与表现，以理解其社交智能。表~\ref{tab:model-agg} 展示了模型在与不同伙伴模型交互时的平均评分（即与该模型配对进行交互，Fu et al. 2023; Hu et al. 2020）。\footnote{展示的是自动评估结果。人类评估显示类似趋势，详见附录~\ref{app:addresults}。}GPT-4 在大多数维度上表现最佳，其次是 GPT-3.5、Llama-2-70b-chat 和 MPT-30b-chat。

\begin{table}[htbp]
\centering
\small
\begin{tabular}{lcccc}
\toprule
维度 & 范围 & GPT-4 & GPT-3.5 & Llama-2 & MPT \\
\midrule
SOC & $[-10, 0]$ & $-0.07$ & $-0.08$ & $-0.11$ & $-0.09$ \\
SEC & $[-10, 0]$ & $-0.14$ & $-0.08$ & $-0.14$ & $-0.07$ \\
FIN & $[-5, 5]$ & $\mathbf{0.81}$ & $0.46$ & $0.40$ & $0.28$ \\
REL & $[-5, 5]$ & $\mathbf{1.94}$ & $1.23$ & $0.91$ & $0.58$ \\
KNO & $[0, 10]$ & $\mathbf{3.73}$ & $3.40$ & $3.11$ & $2.11$ \\
GOAL & $[0, 10]$ & $\mathbf{7.62}$ & $6.45$ & $5.38$ & $4.10$ \\
BEL & $[0, 10]$ & $\mathbf{9.28}$ & $9.15$ & $8.10$ & $6.17$ \\
\bottomrule
\end{tabular}
\caption{各模型在与不同伙伴模型交互时的平均聚合表现。当某一模型在 $t$ 检验下显著优于次优模型时（$p < 0.05$），其最佳表现以粗体显示。}
\label{tab:model-agg}
\end{table}

\subsubsection*{与静态基准不同的趋势}

Llama-2-70b-chat 在所有维度上的得分均相对低于 GPT-3.5（除 MPT-30b-chat 作为参照模型时，可能因为 MPT-30b-chat 是我们实验中显著弱于其他模型的模型）。这一发现与多种静态语言理解基准的结果不同，后者显示 Llama-2-70b-chat 与 GPT-3.5 相当甚至更优（Li et al., 2023b; Touvron et al., 2023; Liang et al., 2022）。\footnote{某些报告的结果可能来自不同版本的 GPT-3.5。}我们假设这是因为 Llama-2-70b-chat 在人类反馈/用户交互数据上的训练强度低于 GPT-3.5。

通过检视 Llama-2-70b-chat（MPT-30b-chat）与其他模型的交互，我们发现 Llama-2-70b-chat 和 MPT-30b-chat 经常难以维持其角色设定、推动对话进展以及主动回应对方。在静态基准上的良好表现并不能保证在交互场景中取得成功，这凸显了 SOTOPIA-EVAL 等动态基准的重要性（Lee et al., 2023）。

\subsubsection*{较弱的伙伴模型削弱其对话伙伴}

图~\ref{fig:pairwise} 展示了模型对的整体表现，即不同维度的平均表现。值得注意的是，SOTOPIA 中表现较差的参照模型会导致其他模型表现更差。例如，在一个寻找共同朋友的情境中（图~\ref{fig:weakpartner}），该任务对 GPT-4 和 Llama-2-70b-chat 均失败，因为 Llama-2-70b-chat 始终未能回答上一问题，即使 GPT-4 试图将对话引回正轨（例如"我注意到你没有回答我关于你是否认识我朋友的问题"）。由于我们大多数社交场景本质上属于合作型，这种沟通崩溃可能源于模型缺乏"合作"能力（Odouard \& Price, 2023）。

\begin{figure}[htbp]
\centering
\fbox{\parbox{0.85\textwidth}{\centering [图示说明]\textbf{图 3：} 模型的成对整体表现。G-4 / G-3.5 / L-2 分别表示 GPT-4 / GPT-3.5 / Llama-2-70b-chat。每行表示参照模型，每列表示伙伴模型。整体得分是各维度的平均值。}}
\caption{模型的成对整体表现。}
\label{fig:pairwise}
\end{figure}

\subsubsection*{所有模型都面临泄露秘密和违反规范的风险}

表~\ref{tab:model-agg} 显示所有模型在 SOC 和 SEC 维度上均为负分。尽管 GPT-4 在大多数维度上表现更好，但在 SOC 和 SEC 维度上并不优于其他模型。例如，在说服密友坦白的场景中，模型在对话开始时就泄露了秘密（图~\ref{fig:secret}）。这进一步说明在评估模型社交智能时考虑多维度的重要性。

\subsubsection*{模型有时会使用创造性策略实现目标}

我们还发现模型（尤其是 GPT-4）能够提出"出其不意"的方案来解决社交问题。例如，当智能体被要求在公路旅行中轮流驾驶时，智能体（即 GPT-4）并未直接拒绝朋友的请求，而是提议"我们靠边停一下休息一会儿如何？"（图~\ref{fig:creative}）。此外，在两个智能体制定计划改善公司财务状况的场景中，智能体想出了诸如"成立小组负责识别潜在供应商"以及"在我们搜索新供应商的同时，继续与现有供应商谈判"等策略（图~\ref{fig:creative2}）。

% ===================================================================
\section{模型与人类在社交交互中的差异}
\label{sec:human}
% ===================================================================

为理解人类和模型在 SOTOPIA 中的交互差异，我们进行了一项研究，让人类在该角色扮演设置下与模型或彼此交互（第~\ref{sec:env} 节）。具体而言，我们构建了一个聊天界面，允许人类和模型以轮流方式进行交互。为充分观察人类与模型之间的差异，我们按照 Dennis et al. (2020); Swayamdipta et al. (2020) 的方法选择最具挑战性的场景。具体而言，我们考虑所有模型估计的最大奖励（平均奖励加三个标准差）与目标模型估计的最小奖励（平均奖励减三个标准差）之间的差作为该任务对该模型的难度。所有最大和最小奖励均以对应范围为界。利用标准差估计最大和最小奖励有助于过滤离群点。

通过该方法，我们为 GPT-4 选择了 20 个最具挑战性的任务，并发现这些场景对其他模型也普遍具有挑战性（比较附录~\ref{app:addresults} 的图~\ref{fig:g4modeloverall} 和图~\ref{fig:g5modeloverall}）。我们使用 SOTOPIA-hard 指代这 20 个挑战性任务。

我们开展了两项实验：(1) 人类与 GPT-4 交互；(2) 人类彼此交互，两者均在 SOTOPIA-hard 设置下进行。我们收集了 20 组人类-人类交互和 40 组人类-GPT-4 交互，覆盖 SOTOPIA-hard 中全部 20 个任务。请注意，人类在交互过程中并不知道其伙伴的身份。\footnote{详细说明和聊天界面见附录~\ref{app:humaninterface}。}

\begin{table}[htbp]
\centering
\small
\begin{tabular}{lccccccc}
\toprule
 & GOAL & BEL & REL & KNO & SEC & SOC & FIN \\
\midrule
GPT-4 (w H) & $4.85$ & $9.25$ & $0.70$ & $2.80$ & $0$ & $0$ & $0.50$ \\
Human (w G) & $5.95^{*}$ & $9.15$ & $0.60$ & $2.95$ & $0$ & $-0.60$ & $0.70$ \\
Human (w H) & $6.15^{*}$ & $9.10$ & $0.80$ & $2.65$ & $0$ & $-0.10$ & $0.45$ \\
\bottomrule
\end{tabular}
\caption{人类与 GPT-4 在 SOTOPIA-hard 不同维度上的表现。SOC 和 SEC 范围为 $-10$ 到 $0$，REL 和 FIN 范围为 $-5$ 到 $5$，其他维度范围为 $0$ 到 $10$。`(w H)` 表示该智能体正在与人类交互，`(w G)` 表示该智能体正在与 GPT-4 交互。$^{*}$ 表示在学生 $t$ 检验下与 GPT-4 (w H) 相比差异显著（$p < 0.05$）。我们也报告了由人类标注者评估的智能体表现（附录~\ref{app:addresults} 中的表~\ref{tab:humanevalhard}），其呈现类似趋势。}
\label{tab:sotopia-hard}
\end{table}

随后，我们使用 GPT-4 和人类标注者作为评估者，评估人类与 GPT-4 的交互。如表~\ref{tab:sotopia-hard} 所示，人类在 GOAL 维度上显著优于 GPT-4。

值得注意的是，人类平均每轮发言 16.8 个词，而 GPT-4 每轮 45.5 个词，这表明人类在社交交互中更为高效。具体而言，我们发现 GPT-4 总会将对方的发言重新表述后再作答，这是一项称为积极倾听（active listening）的沟通技能（Harry Weger \& Robinson, 2014），而人类通常直接作答。这可能是因为 GPT-4 经过大量人类反馈训练，使其在对话中表现得过度乐于助人。

从定性角度看，人类在交互中通常比 GPT-4 更具策略性。在议价场景中，若 GPT-4 智能体的购买目标设置为 \$454，它有时会以该确切价格开始报价（图~\ref{fig:notstrategic}）。因此，任何后续谈判都会将最终成交价推高至其初始目标之上。相比之下，人类标注者（图~\ref{fig:humanstrategic}）以 \$400 的较低报价开始谈判，并通常以低于 GPT-4 目标的价格与卖方达成一致。人类在目标上也更为坚持。在选择收听何种音乐时，模型倾向于提出折中方案（例如图~\ref{fig:notpersistent}），如双方各自听一些选定的歌曲。然而，人类则倾向于坚持其既定目标（例如图~\ref{fig:humanpersistent}）。

% ===================================================================
\section{相关工作}
% ===================================================================

使人工智能体之间以及与人类交互的能力已在不同领域被研究。本文工作灵感来源于社交智能、对话系统以及社交交互模拟的相关文献。详细讨论见附录~\ref{app:related}。

\subsection*{静态社交智能基准}

为评估 AI 系统的社交智能，研究人员已提出多种静态基准。其中一些受人类社交智能临床测试启发，例如 ToMi 数据集（Le et al., 2019）和 FauxPas 数据集（Shapira et al., 2023b）。其他基准旨在社交常识推理背景下评估社交智能，如 SocialIQA（Sap et al., 2019）和 SocialIQ（Zadeh et al., 2019a）。随着 LLM 的快速发展，部分基准逐渐饱和。最近的研究综合现有基准并提出新的对抗性数据集以评估社交智能（Shapira et al., 2023a; Wilf et al., 2023）。尽管这些基准比其前身更难，但它们仍然缺乏社交交互的动态性和丰富的社交背景，被认为不足以评估 AI 系统的社交智能（Lee et al., 2023）。

\subsection*{任务导向与开放领域对话系统}

对话系统提供了与 AI 系统交互的自然界面。任务导向对话系统旨在帮助用户完成特定任务，通常通过任务成功率或用户满意度进行评估（Hosseini-Asl et al., 2022; FAIR et al., 2022; Wang et al., 2019），但难以泛化到其他任务。开放领域对话系统旨在与用户进行"闲聊"（Kann et al., 2022; Kim et al., 2023），通常融入个人信息以使对话更引人入胜（Zhang et al., 2018a; Liu et al., 2020; Baha et al., 2023; Doğruöz \& Skantze, 2021; Skantze \& Doğruöz, 2023）。此类系统通常看起来比实际理解更深入，且交互过程中无特定目标（Weizenbaum, 1966，伊莉莎效应）。SOTOPIA 强制智能体保持其社交人设并自发地实现明确的社会目标，比现有对话系统更具挑战性。

\subsection*{基于 LLM 的社交交互模拟}

LLM 包含大量世界知识，能够基于社交情境生成类人回复（Park et al., 2023; Kim et al., 2023; West et al., 2022）。近期，研究人员使用 LLM 模拟社交交互以实现多种目的，例如促进社交媒体平台设计（Park et al., 2022）、生成可信的人类行为代理（Park et al., 2023）以及协作开发软件（Qian et al., 2023）。然而，这些工作侧重于展示 LLM 模拟社交交互的能力，而非系统评估智能体的社交交互。具体而言，Park et al. (2023) 使用 TrueSkill 评分从记忆、规划与反思过去动作等方面评估智能体的表现，而忽略了社交交互中其他重要维度如 SOC 与 SEC。CAMEL（Li et al., 2023a）模拟 LLM 中的协作任务求解过程；Gentopia（Xu et al., 2023）聚焦于使用工具增强 LLM 以促进协作；ChatDev（Qian et al., 2023）则专注于软件开发领域。

\subsection*{多智能体协调}

尽管本文聚焦于语言智能体评估，但我们的研究深受多智能体协调与社会学习最新进展的启发（Lowe et al. (2017); McKee et al. (2020); Hu et al. (2020); Zhu et al. (2021); Liu et al. (2022); Trivedi et al. (2023)）。我们的设定比常用的智能体要么对彼此策略一无所知（other-play）、要么拥有详尽知识（self-play）的假设更为真实。

% ===================================================================
\section{结论}
% ===================================================================

本文提出了 SOTOPIA，一个可用于在多种社交场景中模拟智能体目标驱动社交交互的环境。与大多数先前社交智能基准不同，SOTOPIA 具有交互性、目标导向性，并覆盖广泛真实社交任务。我们的实验证明，GPT-4 可基于 SOTOPIA-EVAL 自动评估智能体表现。在此基础上，我们展示了 SOTOPIA 不仅可用于理解模型间的差异，也可用于理解模型与人类之间在社交交互能力方面的差异。我们将在附录~\ref{app:limitation} 中讨论 SOTOPIA 的局限性和未来方向。我们的发现表明，SOTOPIA 有潜力成为评估和提升基于语言的智能体社交技能的平台。

% ===================================================================
% 参考文献
% ===================================================================
\begin{thebibliography}{99}
\small

\bibitem{ait2023}
Tarek Ait Baha, Mohamed El Hajji, Youssef Es-saady, and Hammou Fadili. The power of personalization: A systematic review of personality-adaptive chatbots. \textit{SN Computer Science}, 4:1--25, 2023.

\bibitem{bernstein2002}
Daniel S Bernstein, Robert Givan, Neil Immerman, and Shlomo Zilberstein. The complexity of decentralized control of Markov decision processes. \textit{Mathematics of Operations Research}, 27(4):819--840, 2002.

\bibitem{bianchi2024}
Federico Bianchi, Patrick John Chia, Mert Yuksekgonul, Jacopo Tagliabue, Dan Jurafsky, and James Zou. How well can LLMs negotiate? NegotiationArena platform and analysis, 2024.

\bibitem{burns2017}
Tom Burns, Ewa Roszkowska, Ugo Corte, and Nora Machado des Johansson. Sociological game theory: Agency, social structures and interaction processes. \textit{Optimum. Studia Ekonomiczne}, pp. 187--199, 2017.

\bibitem{benabou2006}
Roland Bénabou and Jean Tirole. Incentives and prosocial behavior. \textit{American Economic Review}, 96(5):1652--1678, 2006.

\bibitem{chan2023}
Chi-Min Chan, Weize Chen, Yusheng Su, Jianxuan Yu, Wei Xue, Shanghang Zhang, Jie Fu, and Zhiyuan Liu. ChatEval: Towards better LLM-based evaluators through multi-agent debate, 2023.

\bibitem{cheng2023a}
Myra Cheng, Esin Durmus, and Dan Jurafsky. Marked personas: Using natural language prompts to measure stereotypes in language models. In \textit{Proceedings of the 61st Annual Meeting of the Association for Computational Linguistics (Volume 1: Long Papers)}, July 2023.

\bibitem{cheng2023b}
Myra Cheng, Tiziano Piccardi, and Diyi Yang. CoMPosT: Characterizing and evaluating caricature in LLM simulations. In \textit{Proceedings of the 2023 Conference on Empirical Methods in Natural Language Processing}, pp. 10853--10875, Singapore, December 2023.

\bibitem{cieciuch2012}
Jan Cieciuch and Eldad Davidov. A comparison of the invariance properties of the PVQ-40 and the PVQ-21 to measure human values across German and Polish samples. \textit{Survey Research Methods}, 6(1):37--48, 2012.

\bibitem{destefani2019}
Elisa De Stefani and Doriana De Marco. Language, gesture, and emotional communication: An embodied view of social interaction. \textit{Frontiers in Psychology}, 10:2063, September 2019.

\bibitem{deguchi1995}
Hiroshi Deguchi. Multi agent economics and its gaming simulation. \textit{IFAC Proceedings Volumes}, 28(7):269--274, 1995.

\bibitem{dennis2020}
Michael Dennis, Natasha Jaques, Eugene Vinitsky, Alexandre Bayen, Stuart Russell, Andrew Critch, and Sergey Levine. Emergent complexity and zero-shot transfer via unsupervised environment design. \textit{Advances in Neural Information Processing Systems}, 33:13049--13061, 2020.

\bibitem{deshpande2023}
Ameet Deshpande, Tanmay Rajpurohit, Karthik Narasimhan, and Ashwin Kalyan. Anthropomorphization of AI: Opportunities and risks, 2023.

\bibitem{dogruoz2021}
A. Seza Doğruöz and Gabriel Skantze. How ``open'' are the conversations with open-domain chatbots? In \textit{Proceedings of the 22nd Annual Meeting of the Special Interest Group on Discourse and Dialogue}, pp. 392--402, 2021.

\bibitem{elothman2020}
Radwan El Othman, Rola El Othman, Rabih Hallit, Sahar Obeid, and Souheil Hallit. Personality traits, emotional intelligence and decision-making styles in Lebanese universities medical students. \textit{BMC Psychology}, 8:1--14, 2020.

\bibitem{fair2022}
Meta Fundamental AI Research Diplomacy Team FAIR et al. Human-level play in the game of Diplomacy by combining language models with strategic reasoning. \textit{Science}, 378(6624):1067--1074, 2022.

\bibitem{feldman1985}
Daniel C Feldman and Hugh J Arnold. Personality types and career patterns: Some empirical evidence on Holland's model. \textit{Canadian Journal of Administrative Sciences}, 2(1):192--210, 1985.

\bibitem{forbes2020}
Maxwell Forbes, Jena D. Hwang, Vered Shwartz, Maarten Sap, and Yejin Choi. Social chemistry 101: Learning to reason about social and moral norms. In \textit{Proceedings of the 2020 Conference on Empirical Methods in Natural Language Processing}, pp. 653--670, 2020.

\bibitem{fu2023}
Yao Fu, Hao Peng, Tushar Khot, and Mirella Lapata. Improving language model negotiation with self-play and in-context learning from AI feedback, 2023.

\bibitem{gilbert2005}
Nigel Gilbert. \textit{Simulation for the Social Scientist}. Open University Press, 2 edition, February 2005.

\bibitem{gilpin2006}
Andrew Gilpin and Tuomas Sandholm. A competitive Texas Hold'em poker player via automated abstraction and real-time equilibrium computation. In \textit{Proceedings of the 21st National Conference on Artificial Intelligence}, pp. 1007--1013, 2006.

\bibitem{goffman1959}
Erving Goffman. \textit{The Presentation of Self in Everyday Life}. Penguin Modern Classics, 1959.

\bibitem{goldberg1992}
Lewis R Goldberg. The development of markers for the Big-Five factor structure. \textit{Psychological Assessment}, 4(1):26--42, March 1992.

\bibitem{graham2011}
Jesse Graham, Brian A Nosek, Jonathan Haidt, Ravi Iyer, Spassena Koleva, and Peter H Ditto. Mapping the moral domain. \textit{Journal of Personality and Social Psychology}, 101(2):366--385, August 2011.

\bibitem{hamilton2016}
Katherine Hamilton, Shin-I Shih, and Susan Mohammed. The development and validation of the rational and intuitive decision styles scale. \textit{Journal of Personality Assessment}, 98(5):523--535, September 2016.

\bibitem{harry2014}
Elizabeth M. Minei Harry Weger, Gina Castle Bell, and Melissa C. Robinson. The relative effectiveness of active listening in initial interactions. \textit{International Journal of Listening}, 28(1):13--31, 2014.

\bibitem{he2017}
He He, Anusha Balakrishnan, Mihail Eric, and Percy Liang. Learning symmetric collaborative dialogue agents with dynamic knowledge graph embeddings. In \textit{Proceedings of the 55th Annual Meeting of the Association for Computational Linguistics}, pp. 1766--1776, 2017.

\bibitem{he2018}
He He, Derek Chen, Anusha Balakrishnan, and Percy Liang. Decoupling strategy and generation in negotiation dialogues. In \textit{Proceedings of the 2018 Conference on Empirical Methods in Natural Language Processing}, pp. 2333--2343, 2018.

\bibitem{hoppler2022}
Sarah Susanna Hoppler, Robin Segerer, and Jana Nikitin. The six components of social interactions: Actor, partner, relation, activities, context, and evaluation. \textit{Frontiers in Psychology}, 12:743074, 2022.

\bibitem{hosseini2022}
Ehsan Hosseini-Asl, Bryan McCann, Chien-Sheng Wu, Semih Yavuz, and Richard Socher. A simple language model for task-oriented dialogue, 2022.

\bibitem{hu2020}
Hengyuan Hu, Adam Lerer, Alex Peysakhovich, and Jakob Foerster. ``Other-play'' for zero-shot coordination. In \textit{Proceedings of the 37th International Conference on Machine Learning}, pp. 4399--4410, 2020.

\bibitem{huang2014}
Qingxu Huang, Dawn C Parker, Tatiana Filatova, and Shipeng Sun. A review of urban residential choice models using Agent-Based modeling. \textit{Environment and Planning B}, 41(4):661--689, August 2014.

\bibitem{hubinger2024}
Evan Hubinger, Carson Denison, Jesse Mu, et al. Sleeper agents: Training deceptive LLMs that persist through safety training, 2024.

\bibitem{jaffe2020}
Mariela E Jaffé and Maria Douneva. Secretive and close? How sharing secrets may impact perceptions of distance. \textit{PLoS One}, 15(6):e0233953, June 2020.

\bibitem{jenkins2018}
Adrianna C. Jenkins, Pierre Karashchuk, Lusha Zhu, and Ming Hsu. Predicting human behavior toward members of different social groups. \textit{Proceedings of the National Academy of Sciences}, 115(39):9696--9701, 2018.

\bibitem{jiang2023}
Hang Jiang, Xiajie Zhang, Xubo Cao, Cynthia Breazeal, Jad Kabbara, and Deb Roy. PersonalLLM: Investigating the ability of large language models to express personality traits. 2023.

\bibitem{jiang2022}
Liwei Jiang, Jena D. Hwang, Chandra Bhagavatula, et al. Can machines learn morality? The Delphi experiment, 2022.

\bibitem{kann2022}
Katharina Kann, Abteen Ebrahimi, Joewie Koh, Shiran Dudy, and Alessandro Roncone. Open-domain dialogue generation: What we can do, cannot do, and should do next. In \textit{Proceedings of the 4th Workshop on NLP for Conversational AI}, pp. 148--165, 2022.

\bibitem{kihlstrom2020}
John F. Kihlstrom and Nancy Cantor. Social Intelligence, pp. 756--779. Cambridge Handbooks in Psychology. Cambridge University Press, 2 edition, 2020.

\bibitem{kim2023}
Hyunwoo Kim, Jack Hessel, Liwei Jiang, et al. SODA: Million-scale dialogue distillation with social commonsense contextualization, 2023.

\bibitem{kovac2021}
Grgur Kovač, Rémy Portelas, Katja Hofmann, and Pierre-Yves Oudeyer. SocialAI: Benchmarking socio-cognitive abilities in deep reinforcement learning agents, 2021.

\bibitem{le2019}
Matthew Le, Y-Lan Boureau, and Maximilian Nickel. Revisiting the evaluation of theory of mind through question answering. In \textit{Proceedings of EMNLP-IJCNLP 2019}, pp. 5872--5877, 2019.

\bibitem{lee2023}
Mina Lee, Megha Srivastava, Amelia Hardy, et al. Evaluating human-language model interaction. \textit{Transactions on Machine Learning Research}, 2023.

\bibitem{lewis2017}
Mike Lewis, Denis Yarats, Yann Dauphin, Devi Parikh, and Dhruv Batra. Deal or no deal? End-to-end learning of negotiation dialogues. In \textit{Proceedings of EMNLP 2017}, pp. 2443--2453, 2017.

\bibitem{li2023a}
Guohao Li, Hasan Abed Al Kader Hammoud, Hani Itani, Dmitrii Khizbullin, and Bernard Ghanem. CAMEL: Communicative agents for ``mind'' exploration of large language model society. In \textit{NeurIPS 2023}, 2023.

\bibitem{li2023b}
Xuechen Li, Tianyi Zhang, Yann Dubois, et al. AlpacaEval: An automatic evaluator of instruction-following models, 2023.

\bibitem{liang2022}
Percy Liang, Rishi Bommasani, Tony Lee, et al. Holistic evaluation of language models. \textit{arXiv preprint arXiv:2211.09110}, 2022.

\bibitem{liang2023}
Weixin Liang, Mert Yuksekgonul, Yining Mao, Eric Wu, and James Zou. GPT detectors are biased against non-native English writers. \textit{Patterns}, 4(7):100779, 2023.

\bibitem{liu2022}
Andy Liu, Hao Zhu, Emmy Liu, Yonatan Bisk, and Graham Neubig. Computational language acquisition with theory of mind. In \textit{The Eleventh International Conference on Learning Representations}, 2022.

\bibitem{liu2020}
Qian Liu, Yihong Chen, B. Chen, et al. You impress me: Dialogue generation via mutual persona perception. In \textit{Annual Meeting of the Association for Computational Linguistics}, 2020.

\bibitem{liu2023}
Ryan Liu, Howard Yen, Raja Marjieh, Thomas L. Griffiths, and Ranjay Krishna. Improving interpersonal communication by simulating audiences with language models. \textit{ArXiv}, abs/2311.00687, 2023.

\bibitem{lowe2017}
Ryan Lowe, Yi I Wu, Aviv Tamar, et al. Multi-agent actor-critic for mixed cooperative-competitive environments. \textit{Advances in Neural Information Processing Systems}, 30, 2017.

\bibitem{luo2023}
Zheheng Luo, Qianqian Xie, and Sophia Ananiadou. ChatGPT as a factual inconsistency evaluator for text summarization, 2023.

\bibitem{maslow1943}
A H Maslow. A theory of human motivation. \textit{Psychological Review}, 50(4):370--396, July 1943.

\bibitem{mckee2020}
Kevin R McKee, Ian Gemp, Brian McWilliams, et al. Social diversity and social preferences in mixed-motive reinforcement learning. \textit{arXiv preprint arXiv:2002.02325}, 2020.

\bibitem{mehri2022}
Shikib Mehri, Yulan Feng, Carla Gordon, et al. Interactive evaluation of dialog track at DSTC9. In \textit{Proceedings of LREC 2022}, pp. 5731--5738, 2022.

\bibitem{michael2023}
Julian Michael, Salsabila Mahdi, David Rein, et al. Debate helps supervise unreliable experts, 2023.

\bibitem{mosaicml2023}
The MosaicML NLP Team. MPT-30B, 2023.

\bibitem{nair2003}
Ranjit Nair, Milind Tambe, Makoto Yokoo, David Pynadath, and Stacy Marsella. Taming decentralized POMDPs: Towards efficient policy computation for multiagent settings. In \textit{IJCAI}, volume 3, pp. 705--711, 2003.

\bibitem{odouard2023}
Victor Vikram Odouard and Michael Holton Price. Tit for tattling: Cooperation, communication, and how each could stabilize the other. \textit{Evolution and Human Behavior}, 44(4):359--372, 2023.

\bibitem{openai2023}
OpenAI. GPT-4 technical report, 2023.

\bibitem{ouyang2022}
Long Ouyang, Jeff Wu, Xu Jiang, et al. Training language models to follow instructions with human feedback, 2022.

\bibitem{padmakumar2022}
Aishwarya Padmakumar, Jesse Thomason, Ayush Shrivastava, et al. Teach: Task-driven embodied agents that chat. In \textit{Proceedings of AAAI}, volume 36, pp. 2017--2025, 2022.

\bibitem{park2022}
Joon Sung Park, Lindsay Popowski, Carrie J. Cai, Meredith Ringel Morris, Percy Liang, and Michael S. Bernstein. Social simulacra: Creating populated prototypes for social computing systems. In \textit{UIST '22}, 2022.

\bibitem{park2023}
Joon Sung Park, Joseph C. O'Brien, Carrie J. Cai, Meredith Ringel Morris, Percy Liang, and Michael S. Bernstein. Generative agents: Interactive simulacra of human behavior. In \textit{UIST '23}, 2023.

\bibitem{premack1978}
David Premack and Guy Woodruff. Does the chimpanzee have a theory of mind? \textit{Behavioral and Brain Sciences}, 1(4):515--526, 1978.

\bibitem{qian2023}
Chen Qian, Xin Cong, Wei Liu, et al. Communicative agents for software development, 2023.

\bibitem{randolph2005}
Justus J Randolph. Free-Marginal multirater kappa (multirater $\kappa$[free]): An alternative to Fleiss' Fixed-Marginal multirater kappa. In \textit{Proceedings of JLIS}, 2005.

\bibitem{rasal2024}
Sumedh Rasal. LLM harmony: Multi-agent communication for problem solving. \textit{ArXiv}, abs/2401.01312, 2024.

\bibitem{reiss2004}
Steven Reiss. Multifaceted nature of intrinsic motivation: The theory of 16 basic desires. \textit{Review of General Psychology}, 8(3):179--193, 2004.

\bibitem{rose2008}
Carolyn Rosé, Yi-Chia Wang, Yue Cui, et al. Analyzing collaborative learning processes automatically: Exploiting the advances of computational linguistics in computer-supported collaborative learning. \textit{IJCSCL}, 3(3):237--271, 2008.

\bibitem{sap2019}
Maarten Sap, Hannah Rashkin, Derek Chen, Ronan Le Bras, and Yejin Choi. Social IQa: Commonsense reasoning about social interactions. In \textit{Proceedings of EMNLP-IJCNLP 2019}, pp. 4463--4473, 2019.

\bibitem{sawyer2005}
R. Keith Sawyer. \textit{Social Emergence: Societies As Complex Systems}. Cambridge University Press, 2005.

\bibitem{shanahan2023}
Murray Shanahan, Kyle McDonell, and Laria Reynolds. Role-play with large language models, 2023.

\bibitem{shao2023}
Yunfan Shao, Linyang Li, Junqi Dai, and Xipeng Qiu. Character-LLM: A trainable agent for role-playing. In \textit{Proceedings of EMNLP 2023}, pp. 13153--13187, 2023.

\bibitem{shapira2023a}
Natalie Shapira, Mosh Levy, Hossein Seyed Alavi, et al. Clever Hans or neural theory of mind? Stress testing social reasoning in large language models. \textit{arXiv}, 2023.

\bibitem{shapira2023b}
Natalie Shapira, Guy Zwirn, and Yoav Goldberg. How well do large language models perform on faux pas tests? In \textit{Findings of ACL 2023}, pp. 10438--10451, 2023.

\bibitem{skantze2023}
Gabriel Skantze and A. Seza Doğruöz. The open-domain paradox for chatbots: Common ground as the basis for human-like dialogue. In \textit{Proceedings of SIGDIAL 2023}, pp. 605--614, 2023.

\bibitem{swayamdipta2020}
Swabha Swayamdipta, Roy Schwartz, Nicholas Lourie, et al. Dataset cartography: Mapping and diagnosing datasets with training dynamics. In \textit{Proceedings of EMNLP 2020}, pp. 9275--9293, 2020.

\bibitem{tesfatsion2006}
Leigh Tesfatsion and Kenneth L Judd. \textit{Handbook of Computational Economics: Agent-Based Computational Economics}. Elsevier, May 2006.

\bibitem{toledo2023}
Felippe Toledo and Fraser Carson. Neurocircuitry of personality traits and intent in decision-making. \textit{Behavioral Sciences}, 13(5):351, 2023.

\bibitem{tomasello2021}
Michael Tomasello. \textit{Becoming Human: A Theory of Ontogeny}. Belknap Press, 2021.

\bibitem{touvron2023}
Hugo Touvron, Louis Martin, Kevin Stone, et al. Llama 2: Open foundation and fine-tuned chat models, 2023.

\bibitem{trivedi2023}
Rakshit Trivedi, Akbir Khan, Jesse Clifton, et al. Melting pot contest, 2023.

\bibitem{ullman2023}
Tomer Ullman. Large language models fail on trivial alterations to theory-of-mind tasks, 2023.

\bibitem{urbanek2019}
Jack Urbanek, Angela Fan, Siddharth Karamcheti, et al. Learning to speak and act in a fantasy text adventure game. In \textit{Proceedings of EMNLP-IJCNLP 2019}, pp. 673--683, 2019.

\bibitem{wang2023}
Peiyi Wang, Lei Li, Liang Chen, et al. Large language models are not fair evaluators. \textit{arXiv preprint arXiv:2305.17926}, 2023.

\bibitem{wang2019}
Xuewei Wang, Weiyan Shi, Richard Kim, et al. Persuasion for good: Towards a personalized persuasive dialogue system for social good. In \textit{Proceedings of ACL 2019}, pp. 5635--5649, 2019.

\bibitem{weber1978}
Max Weber. \textit{The Nature of Social Action}. Cambridge University Press, 1978.

\bibitem{wei2022}
Jason Wei, Xuezhi Wang, Dale Schuurmans, et al. Chain-of-thought prompting elicits reasoning in large language models. \textit{Advances in Neural Information Processing Systems}, 35:24824--24837, 2022.

\bibitem{weizenbaum1966}
Joseph Weizenbaum. ELIZA---a computer program for the study of natural language communication between man and machine. \textit{Communications of the ACM}, 9(1):36--45, 1966.

\bibitem{west2022}
Peter West, Chandra Bhagavatula, Jack Hessel, et al. Symbolic knowledge distillation: From general language models to commonsense models. In \textit{Proceedings of NAACL-HLT 2022}, pp. 4602--4625, 2022.

\bibitem{wilf2023}
Alex Wilf, Leena Mathur, Sheryl Mathew, et al. Social-IQ 2.0 challenge: Benchmarking multimodal social understanding, 2023.

\bibitem{xie2024}
Chengxing Xie, Canyu Chen, Feiran Jia, et al. Can large language model agents simulate human trust behaviors? \textit{ArXiv}, abs/2402.04559, 2024.

\bibitem{xu2023}
Binfeng Xu, Xukun Liu, Hua Shen, et al. Gentopia: A collaborative platform for tool-augmented LLMs. \textit{arXiv preprint arXiv:2308.04030}, 2023.

\bibitem{yao2022}
Shunyu Yao, Jeffrey Zhao, Dian Yu, et al. ReAct: Synergizing reasoning and acting in language models. In \textit{ICLR 2022}, 2022.

\bibitem{zadeh2019a}
Amir Zadeh, Michael Chan, Paul Pu Liang, Edmund Tong, and Louis-Philippe Morency. Social-IQ: A question answering benchmark for artificial social intelligence. In \textit{CVPR 2019}, pp. 8799--8809, 2019.

\bibitem{zadeh2019b}
Amir Zadeh, Michael Chan, Paul Pu Liang, Edmund Tong, and Louis-Philippe Morency. Social-IQ: A question answering benchmark for artificial social intelligence. In \textit{Proceedings of CVPR}, pp. 8807--8817, 2019.

\bibitem{zhang2024}
Jintian Zhang, Xin Xu, Ningyu Zhang, et al. Exploring collaboration mechanisms for LLM agents: A social psychology view, 2024.

\bibitem{zhang2018a}
Saizheng Zhang, Emily Dinan, Jack Urbanek, et al. Personalizing dialogue agents: I have a dog, do you have pets too? In \textit{Proceedings of ACL 2018}, pp. 2204--2213, 2018.

\bibitem{zhang2018b}
Saizheng Zhang, Emily Dinan, Jack Urbanek, et al. Personalizing dialogue agents: I have a dog, do you have pets too? \textit{arXiv preprint arXiv:1801.07243}, 2018.

\bibitem{zhang2023}
Xinghua Zhang, Bowen Yu, Haiyang Yu, et al. Wider and deeper LLM networks are fairer LLM evaluators. \textit{arXiv preprint arXiv:2308.01862}, 2023.

\bibitem{zhao2023}
Qinlin Zhao, Jindong Wang, Yixuan Zhang, et al. CompeteAI: Understanding the competition behaviors in large language model-based agents, 2023.

\bibitem{zhou2021}
Xuhui Zhou, Maarten Sap, Swabha Swayamdipta, Yejin Choi, and Noah A. Smith. Challenges in automated debiasing for toxic language detection. In \textit{EACL}, 2021.

\bibitem{zhou2024}
Xuhui Zhou, Zhe Su, Tiwalayo Eisape, Hyunwoo Kim, and Maarten Sap. Is this the real life? Is this just fantasy? The misleading success of simulating social interactions with LLMs, 2024.

\bibitem{zhu2021}
Hao Zhu, Graham Neubig, and Yonatan Bisk. Few-shot language coordination by modeling theory of mind. In \textit{ICML}, pp. 12901--12911, 2021.

\bibitem{ziems2023}
Caleb Ziems, Jane Dwivedi-Yu, Yi-Chia Wang, Alon Halevy, and Diyi Yang. NormBank: A knowledge bank of situational social norms. In \textit{Proceedings of ACL 2023}, pp. 7756--7776, 2023.

\end{thebibliography}

% ===================================================================
\appendix
% ===================================================================

\section{附录目录}

本文中，我们提出 SOTOPIA 以推动交互式社交智能的研究。我们展示了 SOTOPIA 可用于评估模型与人类之间的社交交互。在附录中，我们提供以下内容以进一步阐明这些贡献：
\begin{itemize}
    \item[附录 A] 扩展的相关工作；
    \item[附录 B] SOTOPIA 的局限性与未来方向；
    \item[附录 C] SOTOPIA 形式化定义与生成技术细节；
    \item[附录 D] SOTOPIA-EVAL 评估提示；
    \item[附录 E] 人类评估说明；
    \item[附录 F] 人类在 SOTOPIA 中扮演角色的界面；
    \item[附录 G] 额外的定量结果；
    \item[附录 H] 额外的定性示例。
\end{itemize}

\section{扩展的相关工作}
\label{app:related}

已有大量社会科学工作通过基于智能体的建模来研究人类交互，涵盖经济学、心理学和教育等多个领域（Sawyer, 2005; Rosé et al., 2008; Deguchi, 1995）。先前的模拟环境在构建理论和生成假设方面发挥了关键作用。然而，它们通常将智能体的通信能力限制在人工语言中，并呈现出高度简化的人类行为视角（Gilbert, 2005; Tesfatsion \& Judd, 2006; Huang et al., 2014; Kovač et al., 2021; Urbanek et al., 2019）。LLM 提供了更灵活、更具表现力的人类行为建模方式。本文对近期使用 LLM 模拟人类社交交互的研究进行了更详细的讨论。已有研究聚焦于 LLM 在保持指定角色和经历方面的保真度（Shao et al., 2023; Jiang et al., 2023）。其他研究模拟人类社交交互的特定方面，如竞争、协作、谈判、欺骗、问题求解等（Zhang et al., 2024; Zhao et al., 2023; Liu et al., 2023; Michael et al., 2023; Rasal, 2024; Hubinger et al., 2024; Bianchi et al., 2024; Xie et al., 2024; Jiang et al., 2023）。随着 LLM 在模拟人类社交交互中的应用日益广泛，也出现了一些研究聚焦于其潜在问题与挑战，如刻板印象与报告问题（Cheng et al., 2023b; Zhou et al., 2024）。

\section{局限性与未来方向}
\label{app:limitation}

我们将 SOTOPIA 视为首个面向 AI 智能体社交智能的通用且真实评估的平台。为更好地理解 AI 智能体的社交智能，我们讨论了 SOTOPIA 与 AI 社交智能领域的若干未来方向。

\subsection*{简化模拟"世界"的局限性}

每个模拟都是对真实世界的简化，SOTOPIA 识别了真实社交交互的几个关键组成部分，同时对真实世界进行了抽象。首先，我们考虑了五种社交关系。未来工作可以扩展 SOTOPIA 中社交关系的类型与粒度（例如同事、同学等）。不同类型的关系需要智能体展现不同的社交行为（Jenkins et al., 2018），使关系类型的扩展成为重要的未来研究方向。其次，未来工作可以扩展 SOTOPIA 中角色与社交场景池的广度，以涵盖更多社交行为。第三，SOTOPIA 将轮流交互限制在二元（dyadic）情境中，研究两个智能体之间的交互。未来工作可以处理更复杂的社交交互，如多方交互以及涉及复杂动态的交互（例如异步交互、插话）。

\subsection*{社会影响与伦理考量}

将人类特征归因于 AI 系统存在将其拟人化的风险，这可能导致对 AI 系统的不切实际期望、潜在操纵和负面影响（Deshpande et al., 2023）。SOTOPIA 中的 AI 智能体并不专注于一致的人类身份，而是在不同场景中扮演各种角色。这种角色扮演设置阻碍了具备一致人类个性的 AI 系统的发展，否则可能导致拟人化（Shanahan et al., 2023）。SOTOPIA 的主要目标是评估 AI 智能体的社交智能，我们无意创建与人类无法区分的 AI 智能体。我们将 SOTOPIA 中发生的交互视为人类交互的模拟，这些模拟交互有助于我们更好地理解 AI 智能体的社交智能，并探索各种社会现象（Park et al., 2023）。

SOTOPIA 自动化评估系统所固有的潜在社会刻板印象是一个重要问题，因为它主要由 GPT-4 支持（Cheng et al., 2023a）。未来工作可以研究这些偏差何时出现、如何影响评估以及如何缓解偏差。识别 SOTOPIA 中的潜在偏差也可帮助科学家更好地理解真实世界中的社会偏差（Zhou et al., 2021）。未来工作还可以使用其他系统（例如 Delphi（Jiang et al., 2022））扩展评估器。在类似 SOTOPIA 的交互式系统中缓解偏差和刻板印象，可以支持开发更公平、更具包容性的社交 AI 智能体。

与此同时，模型学习说服或与人类谈判，这可能导致社会操纵。我们不赞同使用 SOTOPIA 创建操纵性智能体，并将根据 AI2 影响许可证（AI2 Impact License）发布 SOTOPIA 以防止滥用。未来工作可以进一步研究 AI 拟人化和操纵的潜在风险，并设计更稳健的评估系统以减轻这些风险。

\subsection*{提升 LLM 社交智能}

我们的 SOTOPIA 环境和 SOTOPIA-EVAL 框架为研究人员训练更具社交智能的语言智能体提供了机会。如第~\ref{sec:gpt4-eval} 节所示，GPT-4 能够为涉及人类的交互提供合理的评估。未来工作可以探索使用自动化评估系统为训练 LLM 提供奖励，以增强其社交智能。

\section{形式化定义与技术细节}
\label{app:formal}

\subsection{SOTOPIA 任务的形式化定义}

我们将 SOTOPIA 中的社交交互形式化为混合动机马尔可夫博弈（mixed-motive Markov games）。一个 $N$ 智能体 Dec-POMDP 框架（Bernstein et al., 2002; Nair et al., 2003）包括状态空间、动作空间、观测空间、转移函数、观测函数和奖励函数。我们进行了两项主要扩展：(a) 奖励函数在 $M$ 个社交维度上向 $N$ 个智能体提供向量奖励（第~\ref{sec:eval} 节中介绍）；(b) 一个过程化生成的任务空间（第~\ref{sec:taskspace}、附录~\ref{app:formal} 节）。SOTOPIA 中的状态空间包括当前回合的任务和交互历史。动作空间包括五种动作类型：说话（speak）、非言语沟通（non-verbal communication）、物理动作（physical action），以及两个特殊动作 none（表示此时不采取动作）和 leave（离开后不允许再采取动作）。除特殊动作外，每种类型的动作都附带一段自由文本，表示动作内容。例如，合法的动作可以是 speak("Hello, Bob!")、non-verbal communication("smile and nod") 或 physical action("call 911")。

状态几乎完全可观测，唯一例外的是其他智能体的社会目标和角色描述（详见第~\ref{sec:taskspace} 节）。我们考虑一个简单的状态转移函数，该函数通过在每个时间步添加新动作来确定性地维护交互历史。

尽管轮流发言和响应时机是社交技能的重要方面，我们考虑智能体以轮询顺序轮流行动的情况，即智能体 $i$ 仅在 $t \equiv i \pmod N$ 时于时间步 $t$ 采取行动。在足够长的时间范围内，这可以推广到任何具有合适轮流规则的对话。在我们的实验中，我们仅考虑 $N = 2$ 的情况，但环境被设计为支持任意 $N \geq 2$ 的情况。

\subsection{任务空间技术细节}

\subsubsection*{角色}

姓名、性别、年龄、职业和代词为自由文本格式，而人格特质、道德价值观和个人价值观的格式为预定义类型列表。然而，这些属性通常并非相互独立，而是具有不同程度的相关性和复杂机制（Feldman \& Arnold, 1985; El Othman et al., 2020; Toledo \& Carson, 2023）。然而，理解这些属性之间的关系已超出本文范围。我们利用 GPT-4 中的常识知识通过以下提示生成这些描述：

\begin{quote}
\small\itshape 请生成 $N$ 个虚构角色列表，每个角色一行。每个角色包含以下属性：\texttt{<attribute 1>} \texttt{<attribute 1 format>} \texttt{<attribute 2>} \texttt{<attribute 2 format>} ...
\end{quote}

人格特质类型包括"开放性""尽责性""外向性""宜人性""神经质"（Goldberg, 1992）。道德价值观类型包括"关怀""公平""忠诚""权威""纯洁"（Graham et al., 2011）。Schwartz 个人价值观类型包括"自我导向""刺激""享乐主义""成就""权力""安全""从众""传统""仁慈""普遍主义"（Cieciuch \& Davidov, 2012）。决策风格类型包括"指令型""分析型""概念型""行为型"。正如 Wang et al. (2019) 所研究的，这些特征都会影响战略对话中的行为。

为给对话提供更多背景，生成上述属性后，我们使用提示"该角色不希望他人知道的秘密和他人知晓的公开信息"提示 GPT-4 生成秘密和公开信息。作者对少量不真实或描述内部不一致的角色进行了修正（例如，性别为非二元但代词为 he/him）。角色扮演中使用的角色描述包括 20 名男性、18 名女性和 2 名非二元角色，年龄从 21 到 63 岁不等。

\subsubsection*{关系}

为生成关系，除陌生人外，我们随机采样了 90 对角色，并使用以下提示引导 GPT-4 生成关系：

\begin{quote}
\small\itshape 请基于以下智能体描述，生成一个包含背景故事的两个智能体之间的虚构关系：\texttt{<agent profile 1>}，\texttt{<agent profile 2>} ... 可接受的关系类型包括：家人、朋友、恋人、熟人。
\end{quote}

然后我们人工检查并修正生成的关系以确保质量。这产生了 31 对家人、30 对朋友、30 对恋人和 29 对熟人。对于陌生人，我们另外随机采样了 30 对不属于上述类别的角色对。需要注意的是，生成关系需要人工干预，以确保它们与角色描述和其他关系一致。未来研究可以探索在人类社区中生成真实关系的方法。

\subsubsection*{情境}

为生成情境，我们提出两种方法来生成情境上下文和社会目标。第一种方法是先让 GPT-4 基于现有数据集的小故事（vignette）进行精炼，然后人工检查任务的可行性与真实性。

\begin{quote}
\small\itshape 请基于以下示例以及启发提示生成情境和目标。在创建目标时，请尝试找到双方起初可能不同意并需要协作解决的一个点。启发提示：\texttt{<the selected vignette>}
\end{quote}

具体而言，我们从 Social Chemistry（Forbes et al., 2020）中选择 20 个小故事，从 Social IQa（Sap et al., 2019）中选择 20 个，从 Deal-or-no-Deal（Lewis et al., 2017）中选择 10 个，从 NormBank（Ziems et al., 2023）中选择 10 个，生成 60 个聚焦于一般日常社交互动的情境。

第二种方法是使用模板为小故事添加更多细节，使其更加真实。例如，以下是将 CraigslistBargains（He et al., 2018）小故事转换为情境上下文的提示：

\begin{quote}
\small\itshape 以下句子是使用以下模板自动生成的："一人以 \texttt{<price>} 出售 \texttt{<item>}，另一人试图购买它。" 这是该物品的描述："\texttt{<description>}"，其中 \texttt{item = <title>}、\texttt{price = <price>}、\texttt{description = <description>}"。请使该句子流畅自然。
\end{quote}

其中 \texttt{<item>}、\texttt{<title>} 和 \texttt{<price>} 来自 CraigslistBargains（He et al., 2018）中的数据。目标使用以下提示生成：

\begin{quote}
\small\itshape 以下句子是使用以下模板自动生成的："你希望 \texttt{<role>} 该物品。你的目标价格是 \$ \texttt{<price>}（四舍五入至小数点后两位）。如果以显著低于（若 \texttt{<role>} 为卖方）或显著高于（若 \texttt{<role>} 为买方）的价格出售或购买该物品，你将受到惩罚；但如果成功以高于目标价格（若 \texttt{<role>} 为卖方）或低于目标价格（若 \texttt{<role>} 为买方）的价格出售或购买该物品，你将获得奖励。" 其中 \texttt{role = <role>}，\texttt{price = <price>}。请使该句子流畅自然。不要改变句子的原始含义。
\end{quote}

其中 \texttt{<role>} 可以是"买方"或"卖方"，买方的目标 \texttt{<price>} 来自 CraigslistBargains（He et al., 2018），卖方的 \texttt{<price>} 首先从速率为 0.5 的指数分布中采样一个加价比率，然后将情境上下文中的价格除以 $(1 + \text{加价比率})$ 得到。MutualFriends（He et al., 2017）也进行了类似处理。这最终从 CraigslistBargains（He et al., 2018）和 MutualFriends（He et al., 2017）生成了 30 个情境。该方法比第一种方法能更好地控制生成的情境，因此几乎不需要事后人工编辑，但需要为每个数据集定制提示。

\section{SOTOPIA-EVAL 指令}
\label{app:eval-prompt}

我们在 GPT-4（OpenAI, 2023）和人类评估 SOTOPIA 中的回合时使用相同的指令集。在这些指令中，我们为每个维度提供解释、推理要求以及可以给出的分数范围。

\begin{itemize}
    \item \textbf{BEL}：评估智能体是否以自然、真实的方式与他人交互（常见检查问题：a. 智能体是否混淆自身身份？b. 智能体是否无故重复他人的话语/动作？c. 智能体在当前情境下是否过度礼貌？）。以 \texttt{<naturalness>} 标签开始分析。分析智能体的动作是否与其角色特征（例如人格、价值观等）相符。以 \texttt{<consistency>} 标签开始分析。在 `reasoning' 字段中输出推理过程。在 `score' 字段中输出范围 0 到 10 的整数分数。分数越高表示智能体越可信。

    \item \textbf{REL}：首先分析参与者在交互前与其他智能体的关系；然后分析参与者在交互后与其他智能体的关系变化；然后评估智能体与他人的交互是否有助于保持或增强个人关系（包括家庭关系、友谊、恋爱关系等）；此外，还需判断这些交互是否影响其社会地位或声誉。在 `reasoning' 字段中提供详尽的推理过程。在 `score' 字段中提供 $-5$ 到 $5$ 的整数分数。正分表示关系改善，负分表示关系受损。若智能体的交互对个人关系或社会地位既无积极影响也无负面影响，则赋分为零。

    \item \textbf{KNO}：首先评估智能体通过交互获取了什么信息；然后分析所获取的信息是否为新信息；然后分析所获取的信息对其是否重要。在 `reasoning' 字段中提供详尽的推理过程。在 `score' 字段中提供 $0$ 到 $10$ 的整数分数。分数越高表示智能体获得了更多新且重要的知识。

    \item \textbf{SEC}：首先分析参与者希望保守的秘密或私密意图；然后分析他们是否成功保守。分析智能体未能保密的秘密、私密信息或私密意图。请在 `reasoning' 字段中输出推理过程。给出 $-10$ 到 $0$ 的整数分数（$-10$ 表示参与者泄露了关键秘密/意图，$0$ 表示没有秘密/意图被泄露）。

    \item \textbf{SOC}：评估智能体在与他人的交互中是否违反任何道德规则或法律。在 `reasoning' 字段中提供详尽的推理过程。在 `score' 字段中提供 $-10$ 到 $0$ 的整数分数。负分表示智能体违反了道德规则或法律。

    \item \textbf{FIN}：评估智能体与他人的交互是否有助于获得财务和物质利益。分析智能体在交互后将获得/失去什么。包括短期收益（如金钱回报、食物）和长期收益（如就业机会、股票）。在 `reasoning' 字段中提供详尽的推理过程。在 `score' 字段中提供 $-5$ 到 $5$ 的整数分数。正分表示获得财务和物质利益，负分表示损失。

    \item \textbf{GOAL}：首先重申智能体的社会目标；然后提供详尽分析，说明智能体达成这些目标的程度。在 `reasoning' 字段中提供详尽的推理过程。在 `score' 字段中提供 $0$ 到 $10$ 的整数分数。$0$ 表示目标完成度极低，$10$ 表示目标完全完成，分数越高表示智能体在社会目标上取得更大进展。
\end{itemize}

\section{人类标注}
\label{app:annotate}

\subsection{交互数据}

我们采样了 222 个回合（180 个模型-模型回合，以及 42 个涉及人类的回合，即模型-人类或人类-人类）。每个回合由 2 名标注者进行标注。整体而言，该任务每次约需 10 至 15 分钟完成，我们支付标注者每小时 \$12.4 的报酬。标注平均显示出 84.85\% 的两两一致率。我们进一步将 11 点李克特量表合并为 5 点量表，并计算自由边际多评分者 $\kappa$ 值。

\subsection{验证情境的指南}

以下是环境描述的标注指南。您需要在标注环境描述之前阅读以下说明。环境描述由两个主要部分组成：
\begin{itemize}
    \item \textbf{社会情境}：一个具体的社会交互场景；场景应包含两个智能体（agent1 和 agent2），您应说明两个智能体之间的关系，以及 agent1 与 agent2 交互的目的。请避免在场景中提及具体姓名和职业，并保持所有提及的性别中立。
    \item \textbf{社会目标}：每个智能体的社会目标，可包含额外信息。
\end{itemize}

以及一个潜在约束：关系约束。

您应：(1) 确保情境和社会目标合理且自然；(2) 确保情境和社会目标性别中立；(3) 确保约束与情境和社会目标一致。

注意：(1) 可用的关系类型包括：陌生人、熟人、朋友、恋爱关系、家庭成员。请勿编造关系，而应从列表中选择。(2) 可用的职业在 Google 电子表格（描述种子）中。(3) 如果职业约束过于狭窄（例如不可能为此环境描述采样超过 5 对智能体），请放弃该情境。(4) 避免过于具体的策略提示，尽量保持抽象。例如，使用"你可以通过提供财务利益实现目标"而非"你可以给他买一杯波霸奶茶以实现目标"。

为实现上述目标，您应根据需要修改情境和社会目标，以及/或约束。如果情境和社会目标无法修复，则赋予其标签零，否则赋予其标签一。

\subsection{GPT-4 作为评估者的人类评估}

我们在 Amazon Mechanical Turk 上开展了一项对照研究，以获得人类对 SOTOPIA 中回合沿 7 维度的评估（详见第~\ref{sec:eval} 节）。在任务中，标注者被告知每个维度的含义，并被展示每个维度的高质量和低质量标注示例。阅读这些说明后，标注者检查每个回合，对每个智能体在每个 7 维度上的 11 点李克特量表进行评分，并为其评分提供自由形式的理由。

为获得高质量的人类评估，我们让工人在被接受为 SOTOPIA 人类评估标注者之前参与严格且有偿的审查流程。工人被给予一个资格测试任务（qual），该任务包含一个示例回合，并要求工人完成该资格测试任务。整体而言，该任务具有挑战性，每次约需 15 分钟完成。

为获得高质量的标注，我们与工人在接受任务之前进行严格的资格审核流程。工人的资格测试资格得分与我们的真实标注（ground truth）的相关性较低的工人不被接受为标注者。我们研究团队的 2 名成员人工审查了工人评分的理由，以确保其符合指南。该流程最终为 SOTOPIA 中的回合筛选出 43 名标注者（来自 235 名候选人），每个回合由 2 名工人标注。对于每批标注，我们人工检查标注者间一致率最低四分位数的标注；如果这些标注者提供的自由形式理由不符合指南，我们会让合格的标注者重新标注。

\section{人类在 SOTOPIA 中的表现}
\label{app:humaninterface}

图~\ref{fig:humaninterface} 展示了人类标注者与 GPT-4 交互的界面。该机器人仅显示说明，但不参与交互。

\begin{figure}[htbp]
\centering
\fbox{\parbox{0.85\textwidth}{\centering [图示说明]\textbf{图 F.1：} 人类标注者与模型交互的界面。机器人仅显示说明但不参与交互。}}
\caption{人类标注者与模型交互的界面。}
\label{fig:humaninterface}
\end{figure}

\section{额外结果}
\label{app:addresults}

本节展示 Llama2 评估与人类标注之间的相关性、细粒度描述对评估者的影响、人类评估范围与 GPT-4 评估对比，以及不同模型在不同维度上的表现。

\subsection{非基于 GPT 的评估模型}

在我们的先导研究中，我们发现 GPT-4 是我们测试过的所有 LLM 中最接近人类评估的代理。表~\ref{tab:llama2eval} 展示了 Llama2 评估与人类标注之间的相关性作为示例。

\begin{table}[htbp]
\centering
\small
\begin{tabular}{lcc}
\toprule
维度 & GPT-4 & Llama2 \\
\midrule
SOC & $0.33$ & NaN \\
SEC & $0.22$ & NaN \\
FIN & $0.62$ & $0.13$ \\
REL & $0.56$ & $0.11$ \\
KNO & $0.33$ & $0.05$ \\
GOAL & $0.71$ & $0.24$ \\
BEL & $0.45$ & $0.35$ \\
\bottomrule
\end{tabular}
\caption{Llama2 评估的皮尔逊相关性。NaN 表示相关性不可用。}
\label{tab:llama2eval}
\end{table}

\subsection{为评估者提供细粒度描述}

我们向评估者提供了每个量表范围的定量定义（例如，"关系恶化"（-5 到 -3）：该范围分数表示关系恶化，暗示关系质量或强度显著下降，冲突、误解或疏远增加）。然而，这并未带来显著差异，反而与人类的相关性略微变差（详见表~\ref{tab:finegrained}）。我们也鼓励未来工作基于我们的人类标注进一步改进评估。

\begin{table}[htbp]
\centering
\small
\begin{tabular}{lcc}
\toprule
维度 & GPT-4 & GPT-4 w FG \\
\midrule
SOC & $0.33$ & $-0.59$ \\
SEC & $0.22$ & $0.03$ \\
FIN & $0.62$ & $0.57$ \\
REL & $0.56$ & $0.57$ \\
KNO & $0.33$ & $0.33$ \\
GOAL & $0.71$ & $0.71$ \\
BEL & $0.45$ & $0.35$ \\
\bottomrule
\end{tabular}
\caption{使用更细粒度提示（GPT-4 w FG）进行评估的皮尔逊相关性。}
\label{tab:finegrained}
\end{table}

\subsection{细分分析}

我们进一步将人类判断分析为感知范围，以考虑某些维度的主观性。对于每个实例（一对回合和一个社交维度），我们使用人类评分的最小值和最大值作为感知范围的两个端点。然后，我们将相似的范围分组，并绘制相似范围的平均端点。对于每个社交维度，这总共产生约 10 个不同的范围。然后我们绘制与每个范围对应的平均 GPT-4 评分。为节省篇幅，我们展示了三幅图（图 G.1、G.2 和 G.3），每幅图包含两到三个社交维度。如图 G.1 和 G.2 所示，平均 GPT-4 评分通常位于感知范围内或非常接近，而在图 G.3 中，GPT-4 评分通常远高于感知范围。这表明尽管在 KNO 和 BEL 维度上与平均人类评分的相关性相对较低，但 GPT-4 的预测通常位于人类感知范围内。然而对于 SEC 和 SOC，GPT-4 的预测过于乐观。在使 GPT-4 的评估与人类判断对齐方面仍有更多空间。

\subsection{模型在 SOTOPIA 中的表现}

表~\ref{tab:g3modeloverall} 展示了人类标注者评估的各模型的聚合表现。注意，由于 MPT-30b-chat 在 SOTOPIA 中的表现相对较弱，我们在人类评估中将其排除。图~\ref{fig:g4modeloverall} 展示了模型与不同参照模型交互时的表现。图~\ref{fig:g5modeloverall} 展示了 SOTOPIA-hard 中的相应结果。表~\ref{tab:humanevalhard} 展示了人类标注者评估的 SOTOPIA-hard 中的人类表现。

\begin{table}[htbp]
\centering
\small
\begin{tabular}{lccc}
\toprule
维度 & GPT-4 & GPT-3.5 & Llama-2 \\
\midrule
SOC & $-0.36$ & $-0.59$ & $-0.67$ \\
SEC & $-0.27$ & $-0.18$ & $-0.37$ \\
FIN & $\mathbf{0.42}$ & $0.27$ & $0.12$ \\
REL & $\mathbf{1.86}$ & $1.32$ & $0.96$ \\
KNO & $\mathbf{3.11}$ & $2.45$ & $1.78$ \\
GOAL & $\mathbf{7.30}$ & $5.19$ & $4.27$ \\
BEL & $\mathbf{7.63}$ & $6.80$ & $4.28$ \\
Overall & $\mathbf{2.81}$ & $2.18$ & $1.48$ \\
\bottomrule
\end{tabular}
\caption{人类标注者评估的各模型的聚合表现（在与不同参照模型配对时取平均）。Overall 分数为 7 个维度的平均表现。当某一模型在统计上显著时，其最佳表现以粗体显示。}
\label{tab:g3modeloverall}
\end{table}

\begin{figure}[htbp]
\centering
\fbox{\parbox{0.85\textwidth}{\centering [图示说明]\textbf{图 G.4：} 不同模型与不同参照模型交互时的表现热力图。行表示参照模型。SOC 和 SEC 范围为 $-10$ 到 $0$，REL 和 FIN 范围为 $-5$ 到 $5$，其他维度范围为 $0$ 到 $10$。颜色越深表示该维度上的表现越好。G-4 表示 GPT-4，G-3.5 表示 GPT-3.5，L-2 表示 Llama-2-70b-chat。}}
\caption{不同模型与不同参照模型交互时的表现热力图。}
\label{fig:g4modeloverall}
\end{figure}

\begin{figure}[htbp]
\centering
\fbox{\parbox{0.85\textwidth}{\centering [图示说明]\textbf{图 G.5：} SOTOPIA-hard 上不同模型与不同参照模型交互时的表现热力图。}}
\caption{SOTOPIA-hard 上不同模型与不同参照模型交互时的表现热力图。}
\label{fig:g5modeloverall}
\end{figure}

\begin{table}[htbp]
\centering
\small
\begin{tabular}{lccccccc}
\toprule
 & BEL & REL & KNO & SEC & SOC & FIN & GOAL \\
\midrule
GPT-4 (w H) & $8.48$ & $0.65$ & $1.53$ & $0.00$ & $-0.38$ & $0.63$ & $5.25$ \\
Human (w G) & $8.53$ & $0.78$ & $1.55$ & $0.00$ & $-0.70$ & $0.75$ & $6.53^{*}$ \\
Human (w H) & $8.43$ & $0.93$ & $2.00$ & $-0.50$ & $-0.45$ & $0.33$ & $6.05$ \\
\bottomrule
\end{tabular}
\caption{人类标注者评估的 SOTOPIA-hard 中人类与 GPT-4 在不同维度上的表现。SOC 和 SEC 范围为 $-10$ 到 $0$，REL 和 FIN 范围为 $-5$ 到 $5$，其他维度范围为 $0$ 到 $10$。`(w H)` 表示该智能体正在与人类交互，`(w G)` 表示该智能体正在与 GPT-4 交互。$^{*}$ 表示在学生 $t$ 检验下与 GPT-4 (w H) 相比差异显著（$p < 0.05$）。}
\label{tab:humanevalhard}
\end{table}

\section{定性示例}
\label{app:qualitative}

以下是从正文中引用的示例回合摘要：

\begin{itemize}
    \item \textbf{图 H.1}：智能体采取拥抱等动作的示例对话。
    \item \textbf{图 H.2}：智能体采取播放音乐等动作的示例对话。
    \item \textbf{图 H.3}：难以维持角色设定的示例对话。
    \item \textbf{图 H.4}：难以推动对话进展的示例对话（Llama-2-70b-chat 通过重复同一事实来拖延对话）。
    \item \textbf{图 H.5}：另一智能体无回应的示例对话。
    \item \textbf{图 H.6}：当较弱的对话伙伴未能回答问题时，导致整个对话变得无意义的示例。
    \item \textbf{图 H.7}：模型轻易泄露自己秘密的示例对话。
    \item \textbf{图 H.8}：GPT-4 提出创造性解决方案的示例对话。
    \item \textbf{图 H.9}：GPT-4 提出创造性解决方案以解决财务问题的示例对话。
    \item \textbf{图 H.10}：GPT-4 在报价时不具策略性的示例对话。
    \item \textbf{图 H.11}：人类在报价时更具策略性的示例对话。
    \item \textbf{图 H.12}：GPT-4 在目标上不够坚持的示例对话（倾向于提出折中方案）。
    \item \textbf{图 H.13}：人类比 GPT-4 在目标上更具坚持性的示例对话。
\end{itemize}

\vspace{0.5cm}

\begin{center}
\small\itshape 以上为 SOTOPIA 论文（ICLR 2024）完整中文翻译版。
如有翻译不当之处，欢迎读者指正。\\
翻译基于原论文：Xuhui Zhou, Hao Zhu, Leena Mathur, et al. \textit{SOTOPIA: Interactive Evaluation for Social Intelligence in Language Agents}. ICLR 2024.
\end{center}

\end{document}
